\documentclass[11pt]{article}
\usepackage[margin=1in]{geometry}
\usepackage{amsmath,amssymb,mathtools}
\usepackage{booktabs}
\usepackage{enumitem}
\usepackage{microtype}
\usepackage{xcolor}
\usepackage[colorlinks=true,linkcolor=blue!50!black,citecolor=blue!50!black,urlcolor=blue!50!black]{hyperref}
\usepackage{parskip}

\newcommand{\E}{\mathbb{E}}
\newcommand{\Prob}{\mathbb{P}}
\newcommand{\ind}[1]{\mathbb{1}_{\{#1\}}}
\newcommand{\sgn}{\operatorname{sgn}}
\newcommand{\Rank}{\operatorname{CsRank}}
\newcommand{\TsRank}{\operatorname{TsRank}}
\DeclareMathOperator{\EMA}{EMA}
\DeclareMathOperator{\VWAP}{VWAP}
\DeclareMathOperator{\Std}{Std}
\DeclareMathOperator{\Corr}{Corr}
\DeclareMathOperator{\Cov}{Cov}
\DeclareMathOperator{\Mean}{Mean}
\DeclareMathOperator{\Skew}{Skew}
\DeclareMathOperator{\Kurt}{Kurt}

% Simple algorithm environment (prose steps, no code)
\newcounter{algo}
\newenvironment{algo}[1]{%
  \refstepcounter{algo}%
  \par\medskip\noindent\textbf{Algorithm \thealgo\ (#1).}%
  \begin{enumerate}[label=\arabic*.,leftmargin=2em,itemsep=1pt,topsep=3pt]}%
  {\end{enumerate}\medskip}

% Factor heading
\newcounter{fac}
\newcommand{\factor}[2]{%
  \refstepcounter{fac}%
  \par\medskip\noindent\textbf{F\thefac\ --- #1} \hfill {\small[#2]}\par\nopagebreak}

\title{\textbf{Alpha-Prediction Factors for Corporate Bonds}\\[4pt]
\large A compendium of signals extracted from limit-order-book microstructure\\ and formulaic alpha-mining literature, translated to an OTC bond tape}
\author{Compiled factor reference}
\date{August 2026}

\begin{document}
\maketitle

\begin{abstract}
\noindent This document extracts the return-predictive (``alpha'') content of fifteen papers spanning limit-order-book (LOB) microstructure, order-flow point-process modeling, price impact, flow toxicity, sparse-orderbook electricity forecasting, and machine-generated formulaic alphas, and restates each signal in a form computable on a corporate-bond data set (a TRACE-style trade tape plus dealer composite quotes). Scope is restricted to \emph{prediction of future price or spread changes}; quote-placement, execution-scheduling, and market-making control problems are deliberately excluded, though predictive state variables arising inside those frameworks are retained. Each factor is given with (i) a mathematical definition, (ii) the economic intuition, and (iii) the bond-market translation. Short, language-agnostic reproduction algorithms follow in Section~\ref{sec:algos}, and a source-mapping table closes the document.
\end{abstract}

\tableofcontents

%======================================================================
\section{Data model, clocks, and notation}\label{sec:notation}

\paragraph{Instruments and tape.} Fix a bond universe $\mathcal{B}$. For bond $b$ the trade tape is a marked point process $\{(t_k, P_k, V_k, d_k)\}_k$ with execution time $t_k$, clean price $P_k$, par volume $V_k$, and direction flag
\[
d_k \in \{+1 \ (\text{customer buy}),\ -1 \ (\text{customer sell}),\ 0 \ (\text{interdealer})\},
\]
as reported on a TRACE-style feed. Where available, a composite quote process supplies bid $P^B_t$, ask $P^A_t$, and contributed sizes $q^B_t, q^A_t$ (e.g.\ aggregated dealer axes or CBBT/CP+-style composites). The mid is $m_t = \tfrac12(P^B_t + P^A_t)$ and the quoted spread $s_t = P^A_t - P^B_t$; when quotes are unavailable, replace $m_t$ by a robust trade-based fair value (F15--F18) and $s_t$ by the realized effective spread (F22).

\paragraph{Returns.} Simple returns $r_{t,\delta} = P_{t+\delta}/P_t - 1$ on price, or spread changes $\Delta s^{OAS}$ where spread-space prediction is preferred. All factors below are stated in price space; a sign flip converts to spread space.

\paragraph{Clocks.} Three clocks aggregate the tape (following the toxic-flow and VPIN literature). This choice is \emph{essential} for bonds, where calendar time is mostly empty:
\begin{itemize}[itemsep=1pt]
\item \textbf{Time clock:} window $=$ last $\delta$ units of calendar time.
\item \textbf{Volume clock:} window $=$ most recent prints totalling $V^\ast$ of par volume.
\item \textbf{Transaction clock:} window $=$ last $k$ prints.
\end{itemize}
A \emph{fallback-window rule} (from the sparse-orderbook electricity literature) handles empty windows: given nested horizons $\delta_1 < \delta_2 < \dots < \delta_W$, compute each feature on the shortest window containing at least one print; discard the observation only if the full-history window is empty.

\paragraph{Operator grammar.} Cross-sectional and time-series operators from the alpha-mining papers, used throughout:
\begin{center}
\small
\begin{tabular}{@{}lll@{}}
\toprule
Operator & Type & Meaning \\
\midrule
$\operatorname{Ref}(x,t)$, $\Delta(x,t) = x - \operatorname{Ref}(x,t)$ & TS & lag; $t$-period difference \\
$\Mean(x,t)$, $\operatorname{Med}(x,t)$, $\operatorname{Sum}(x,t)$ & TS & rolling mean / median / sum \\
$\Std(x,t)$, $\operatorname{Var}(x,t)$, $\operatorname{Mad}(x,t)$ & TS & rolling dispersion; $\operatorname{Mad}=\E|x-\E x|$ \\
$\operatorname{Min}(x,t)$, $\operatorname{Max}(x,t)$ & TS & rolling extremes \\
$\operatorname{WMA}(x,t)$, $\EMA(x,t)$ & TS & weighted / exponential moving average \\
$\Corr(x,y,t)$, $\Cov(x,y,t)$ & TS & rolling correlation / covariance \\
$\TsRank(x,t)$ & TS & rank of current value within its own last $t$ values \\
$\Rank(x)$ & CS & cross-sectional rank across bonds on a date \\
$\operatorname{Greater}$, $\operatorname{Less}$, $\operatorname{IfElse}$, $\operatorname{Abs}$, $\operatorname{Log}$, $\sgn$ & CS & elementwise \\
\bottomrule
\end{tabular}
\end{center}
Cross-sectional ranking should be performed \emph{within homogeneous buckets} (rating $\times$ maturity $\times$ sector, or versus issuer curve) so that ranks compare like with like.

%======================================================================
\section{Signed-flow and trade-sign factors}\label{sec:flow}

\factor{Last-trade sign}{Muni Toke--Yoshida}
\[
\epsilon_t = d_{k(t)}, \qquad k(t) = \max\{k : t_k \le t,\ d_k \neq 0\}.
\]
\emph{Intuition.} Order flow is autocorrelated: the sign of the most recent customer trade predicts the sign of the next one, and hence the direction of the next price adjustment. In the intensity-ratio study this covariate is never dropped by model selection.
\emph{Bond translation.} Directly computable from the TRACE buy/sell flag. Because bonds trade sparsely, decay the signal with time-since-last-trade (see F3) rather than using the raw sign indefinitely.

\factor{Signed-spread interaction}{Muni Toke--Yoshida}
\[
X_t = s_t \,\epsilon_t
\qquad\text{or, in regime form,}\qquad
X_t = \ind{s_t > \bar s}\,\epsilon_t ,
\]
with $\bar s$ a trailing mean of the spread.
\emph{Intuition.} The empirical finding: when the spread is wide, book-imbalance signals flatten and the \emph{persistence of the last trade sign strengthens}. The interaction term, not the spread alone, carries the information; model selection retains $(i, \epsilon, s\epsilon)$ and essentially never $(i,\epsilon,s)$.
\emph{Bond translation.} Bonds are a permanently wide-spread market, so sign-persistence terms should be weighted \emph{up} relative to imbalance terms; the interaction lets the data set the mix bond-by-bond.

\factor{Exponentially decayed signed flow (flow pressure)}{Hawkes covariates; Bank--Cartea--K\"orber}
\[
F_t(\beta) \;=\; \sum_{t_k \le t} e^{-\beta\,(t - t_k)}\; d_k\, V_k ,
\]
computed for a small grid of decay rates $\beta \in \{\beta_{\text{fast}}, \beta_{\text{med}}, \beta_{\text{slow}}\}$ (e.g.\ half-lives of 1 day, 1 week, 1 month for bonds).
\emph{Intuition.} A tradable scalar summary of the mutually exciting order-flow intensities: recent net customer buying predicts continued buying and upward price adjustment at short horizons, and reversion at long horizons once the flow is inventory-driven (see F13).
\emph{Bond translation.} Use signed par volume from TRACE; normalize by amount outstanding or by trailing average volume to make $F_t$ comparable across bonds. The pair $(F_t(\beta_{\text{fast}}), F_t(\beta_{\text{slow}}))$ separates continuation from reversion horizons.

\factor{Per-type Hawkes intensity covariates}{Muni Toke--Yoshida}
For each flow type $i$ (customer buy, customer sell, interdealer; optionally split aggressive vs.\ non-aggressive by whether the print is at/through the far quote),
\[
H_i(t) \;=\; \log\!\Big( \mu_i + \textstyle\int_0^t \alpha_i\, e^{-\beta_i (t-u)}\, dN^i_u \Big),
\]
with $(\mu_i,\alpha_i,\beta_i)$ fitted by maximum likelihood per bond or pooled per liquidity bucket.
\emph{Intuition.} The fitted self-/cross-excitation intensity of each order type is itself a covariate: elevated buy-side intensity relative to sell-side intensity predicts the side of the next trade. In the marked-point-process study, the covariate sets $(Z_1, H_{\text{aggr.buy}}, H_{\text{aggr.sell}}, H_{\text{buy}}, H_{\text{sell}})$ dominate model selection.
\emph{Bond translation.} Fit on the TRACE arrival times per side. For very sparse bonds, pool the Hawkes parameters within issuer or rating bucket and keep only the state $H_i(t)$ bond-specific.

\factor{Buy/sell separated print streams and their gap}{OrderFusion}
Maintain the two sides as separate series and use both their difference and their interaction:
\[
\VWAP^{+}_{[t-\delta,t]} = \frac{\sum_{k:\,d_k=+1} P_k V_k}{\sum_{k:\,d_k=+1} V_k},
\qquad
\VWAP^{-}_{[t-\delta,t]} \ \text{analogously},
\]
\begin{gather*}
G_t = \VWAP^{+} - \VWAP^{-} \quad (\text{realized effective spread; see F22}),\\
\Pi_t = \frac{\sum_{d_k=+1} V_k - \sum_{d_k=-1} V_k}{\sum_{d_k\neq 0} V_k} \quad (\text{volume pressure}).
\end{gather*}
\emph{Intuition.} Aggregating the two sides into one series destroys the buy--sell interaction structure, which is where much of the predictive content lives; the electricity work encodes the sides as two separate 2-D streams for exactly this reason.
\emph{Bond translation.} TRACE side flags make this free. $\Pi_t$ is the trade-based analogue of book imbalance; $G_t$ doubles as a liquidity factor.

%======================================================================
\section{Imbalance and quote-based factors}\label{sec:imb}

\factor{Best-quote imbalance}{Lehalle--Neuman; Muni Toke--Yoshida $Z_1$; Palguna--Pollak SI}
\[
I_t = \frac{q^B_t - q^A_t}{q^B_t + q^A_t} \in [-1,1].
\]
\emph{Intuition.} Standing bid-side size in excess of ask-side size predicts an upward next move; empirically the signal behaves like a mean-reverting Ornstein--Uhlenbeck process,
\[
dI_t = -\gamma I_t\,dt + \sigma_I\, dW_t,
\]
so its value decays at a known rate --- important when the prediction horizon is longer than the signal half-life $\log 2/\gamma$.
\emph{Bond translation.} Replace queue sizes with (i) counts/sizes of dealer axes bid-wanted vs.\ offered, (ii) direction imbalance of recent inquiry flow, or (iii) size asymmetry among composite-quote contributors. Weaker than in a firm-quote LOB (bond quotes are indicative), but the OU decay structure carries over and calibrates the usable horizon.

\factor{Multi-level, multi-lookback book-pressure grid}{Li et al.}
\[
\operatorname{OBP}(n,\ell)_t \;=\;
\frac{\sum_{\tau=0}^{n}\sum_{j=1}^{\ell} \text{BidSize}_{j,\,t-\tau}}
     {\sum_{\tau=0}^{n}\sum_{i=1}^{\ell} \text{AskSize}_{i,\,t-\tau}},
\]
computed on a grid of lookbacks $n$ and depth levels $\ell$ and fed \emph{as a vector} to a classifier (no single calibration).
\emph{Intuition.} Persistent one-sided pressure across depth levels and across time predicts the direction of the mid; letting the learner weight the $(n,\ell)$ grid avoids committing to one parameterization.
\emph{Bond translation.} Replace ``levels'' by \emph{dealers}: aggregate axe/quote size across contributing dealers, over lookbacks of $\{1,5,21\}$ days. The grid-not-scalar design principle transfers unchanged.

\factor{Order-flow imbalance (event-signed quote changes)}{Palguna--Pollak, after Cont et al.}
At quote-update events $T_i$,
\[
e_{i+1} = \ind{P^B_{i+1} \ge P^B_i}\, q^B_{i+1} - \ind{P^B_{i+1} \le P^B_i}\, q^B_i
        - \ind{P^A_{i+1} \le P^A_i}\, q^A_{i+1} + \ind{P^A_{i+1} \ge P^A_i}\, q^A_i,
\qquad
\operatorname{OFI}_t = \sum_{i=t-E+1}^{t} e_i .
\]
\emph{Intuition.} Signed changes in posted supply and demand --- not just standing levels --- measure the net arrival of buying vs.\ selling interest. Predicts sign of the next mid move.
\emph{Bond translation.} Apply to the \emph{composite} bid/ask: a dealer lifting their contributed bid is information even when no trade prints. Event-clocked by construction, hence robust to sparse activity.

%======================================================================
\section{Reversal and price-pressure factors}\label{sec:rev}

\factor{Moving-average deviation reversal}{FactorMiner 002/006}
\[
X_t = -\,\frac{P_t - \EMA(P,\,n)_t}{\EMA(P,\,n)_t}
\qquad\text{and}\qquad
X_t = -\,\frac{P_t - \VWAP_{[t-\delta,t]}}{\VWAP_{[t-\delta,t]}} .
\]
\emph{Intuition.} Price above its own recent average reverts. In credit, the mechanism is concrete: prints away from fair value are dealer-intermediated liquidity events whose price concession is subsequently recouped.
\emph{Bond translation.} Use trade-print price vs.\ VWAP of recent TRACE prints (volume clock recommended); horizon days-to-weeks. Among the most robust documented corporate-bond alphas.

\factor{Volume-conditioned reversal}{FactorMiner 004/008}
\begin{gather*}
X_t = \Rank\!\Big(\frac{V_t}{\Mean(V,\,n)_t}\Big)\cdot \Rank(-\,r_t)
\qquad\text{or}\\
X_t = -\,\Rank\!\Big(\frac{r_t}{\Std(r,\,n)_t + \varepsilon}\Big)\cdot \Rank\!\Big(\frac{V_t}{\Mean(V,\,n)_t}\Big).
\end{gather*}
\emph{Intuition.} Reversal is strongest when the move was accompanied by abnormal volume: high-volume moves are predominantly liquidity demand, and the liquidity provider's compensation is the subsequent reversion. The second form standardizes the return by its own volatility before ranking.
\emph{Bond translation.} TRACE volume spike $\times$ price drop $\Rightarrow$ buy. Normalize volume by amount outstanding. This is the single highest-conviction transfer in the document.

\factor{Range-position (stochastic-oscillator) reversal}{FactorMiner 001/005/018}
\[
K_t(n) = \frac{P_t - \operatorname{Min}(P,\,n)_t}{\operatorname{Max}(P,\,n)_t - \operatorname{Min}(P,\,n)_t + \varepsilon} \in [0,1],
\qquad X_t = -\,\Rank\big(K_t(n)\big),
\]
optionally interacted with $\TsRank$ of volume or of $\Cov(r, V, n)$ as in the mined variants.
\emph{Intuition.} Position within the recent trading range proxies over-extension; the top of the range predicts negative drift and vice versa.
\emph{Bond translation.} Computable on weekly bars where bond sparsity forces coarse sampling; combine with F16 extremes.

\factor{Conditional reversal/momentum switch}{FactorMiner 019/031/034}
\[
X_t = \operatorname{IfElse}\!\Big( |{\textstyle\frac{P_t - \VWAP_t}{\VWAP_t}}| > c,\;
 -\Rank\big(P_t - \operatorname{Max}(P,n)_t\big),\;
 -\Rank\big(P_t - \operatorname{Min}(P,n)_t\big)\Big).
\]
\emph{Intuition.} When price is dislocated from VWAP beyond a threshold, fade the extreme; otherwise trade the drift toward the range boundary. Encodes a regime switch between pressure-reversion and drift regimes.
\emph{Bond translation.} Threshold $c$ in spread terms (e.g.\ a fraction of typical bid--offer) rather than a fixed percentage.

\factor{Impact-state de-pressured value (propagator residual)}{Webster / Obizhaeva--Wang}
Model the transient impact state per bond,
\[
dI_t = -\beta\, I_t\, dt + \lambda\, dQ_t, \qquad dQ_t = d_k V_k \ \text{at print times},
\]
i.e.\ $I_t = \lambda \sum_{t_k \le t} e^{-\beta (t-t_k)} d_k V_k$ (note $I_t = \lambda F_t(\beta)$ from F3), and define the de-pressured value and the factor
\[
\widehat{P}_t = P_t - I_t, \qquad X_t = -\,I_t .
\]
\emph{Intuition.} Observed price $=$ fundamental $+$ transient impact; the transient part decays at rate $\beta$, so the expected return over horizon $h$ is $-I_t\,(1-e^{-\beta h})$. This is the model-based twin of F9: reversal \emph{is} impact decay. The key structural result is myopia --- for a large model class the mapping from flow to expected reversion is a simple function of the impact level and its decay rate, so no dynamic optimization is needed to use it as a predictor.
\emph{Bond translation.} $\beta$ is slow for bonds (impact decays over days, not minutes); estimate $(\lambda,\beta)$ per liquidity bucket by regressing future returns on decayed signed flow at several horizons.

\factor{Price--volume rank divergence}{FactorMiner 003/020/032; AlphaGen operators}
\begin{gather*}
X_t = \Rank\big(\Delta(V,\,k)_t\big) - \Rank\Big(\frac{P_t - \VWAP_t}{\VWAP_t}\Big),\\
X_t = -\,\TsRank\Big( \Rank\big(\Delta(r,\,2)\big) - \Rank\big(\Delta(V,\,2)\big),\, n \Big),
\end{gather*}
\[
X_t = -\,\TsRank\!\Big( \Cov\big(\TsRank(P,\,n),\, \TsRank(V,\,n),\, k\big),\, n\Big)\cdot \TsRank\!\big(K_t(n),\, n\big).
\]
\emph{Intuition.} Volume rising without price confirmation flags accumulation/distribution; sustained co-movement of price and volume ranks near the top of the range flags exhaustion.
\emph{Bond translation.} Weekly TRACE aggregates; untested in credit but computationally free once F9--F11 are built.

%======================================================================
\section{Fair-value statistics on a sparse tape}\label{sec:fv}

The sparse-orderbook electricity studies extract price/volume statistics on both sides over nested lookbacks $\delta_w \in \{1,5,15,60,180,\infty\}$ (minutes there; scale to hours/days for bonds) and select with an $\ell_1$ penalty. The headline empirical result transfers as a prior: \emph{price percentiles ($\approx 30\%$ of selected importance), minima ($\approx 26\%$) and maxima ($\approx 22\%$) dominate; volume statistics contribute little; and medium lookbacks beat the shortest window}, because the freshest print is noisy.

\factor{Trade-price percentiles}{orderbook feature learning}
\[
Q^{(p)}_{t}(\delta) = \operatorname{percentile}_{p}\{P_k : t_k \in [t-\delta, t]\},
\qquad p \in \{10, 25, 45, 50, 55, 75, 90\}\%,
\]
computed separately for customer-buy, customer-sell, and all prints. Factors: the level $Q^{(50)}$ as fair value; the deviation $(P_t - Q^{(50)}_t)/Q^{(50)}_t$ as a reversal signal; the inter-percentile width $Q^{(75)}-Q^{(25)}$ as a dispersion/liquidity state.
\emph{Bond translation.} The median of recent TRACE prints is a robust matrix-price-free fair value; percentile deviations are outlier-resistant versions of F9.

\factor{Rolling extremes and their spread}{orderbook feature learning; FactorMiner}
\[
\operatorname{Min}(P,\delta)_t,\quad \operatorname{Max}(P,\delta)_t,\quad
R_t(\delta) = \operatorname{Max} - \operatorname{Min}.
\]
Used directly (support/resistance content), inside $K_t$ (F11), and as the volatility proxy $R_t/Q^{(50)}_t$.

\factor{Nested-window VWAPs and last/first prices}{OrderFusion; orderbook feature learning}
\[
\VWAP_{[t-\delta_w,\,t]} \ \ \text{for } \delta_w \in \{\delta_1,\dots,\delta_W\},\qquad
P^{\text{last}}_t(\delta_w),\ P^{\text{first}}_t(\delta_w).
\]
\emph{Intuition.} The trade-weighted average over a medium window is the strongest single fair-value estimate in the sparse-market studies; the last print alone is weakly efficient but noisy. The \emph{vector} across windows, with the fallback rule of Section~\ref{sec:notation}, lets a sparse learner choose the horizon.
\emph{Bond translation.} Windows of $\{1, 5, 21, 63\}$ trading days; per side and combined.

\factor{Liquid-to-illiquid signal propagation}{orderbook feature learning (asymmetric generalization)}
For an illiquid bond $b$ of issuer $j$, let $L(j)$ be the issuer's most liquid bond (or the on-the-run index constituent). Define
\[
X^{(b)}_t = X^{(L(j))}_t \quad \text{for any factor above,}
\qquad\text{or the basis } \ Z_t = \tilde P^{(b)}_t - \big(\text{curve-implied price from } L(j)\big).
\]
\emph{Intuition.} The empirical asymmetry: feature sets and models estimated on the liquid market transfer to the illiquid one, \emph{but not the reverse}. Hence: estimate on liquid bonds, propagate signals along the issuer curve, and trade the illiquid sibling's convergence to the liquid-implied value.

%======================================================================
\section{Liquidity, illiquidity, and resilience factors}\label{sec:liq}

\factor{Amihud illiquidity}{FactorMiner 013; standard credit literature}
\[
\operatorname{ILLIQ}_t(n) = \Mean\!\Big( \frac{|r_k|}{V_k},\, n\Big)_t .
\]
\emph{Intuition.} Price move per unit of volume; a priced characteristic (liquidity premium) and a conditioning variable for every flow factor above (impact $\lambda$ in F13 is proportional to it).

\factor{Resilience (illiquidity improvement)}{FactorMiner 014}
\[
X_t = \Delta_k\!\left( \frac{ \Mean\big( |r|/V,\, n\big)_t }{\, |r_t|/V_t + \varepsilon\,} \right),
\]
the $k$-period change in the ratio of trailing-average to current illiquidity.
\emph{Intuition.} Not the level but the \emph{dynamics}: a bond whose price-per-volume sensitivity is improving relative to its own history is absorbing flow better; deteriorating resilience predicts spread widening.

\factor{Amount-efficiency family}{FactorMiner 073--099}
\[
\operatorname{AmtEff}_t = \frac{|r_{t}|}{\text{Turnover}_t + \varepsilon},
\qquad \text{velocity} = \Delta_k \operatorname{AmtEff},
\qquad \text{acceleration} = \Delta_k^2 \operatorname{AmtEff},
\]
with turnover $=$ volume / amount outstanding. Refinements of F19 that the mining agent rediscovered repeatedly; use one or two, orthogonalized against F19.

\factor{Realized effective spread}{OrderFusion side separation; TRACE literature}
\[
S^{\mathrm{eff}}_t(\delta) = \VWAP^{+}_{[t-\delta,t]} - \VWAP^{-}_{[t-\delta,t]},
\]
customer-buy VWAP minus customer-sell VWAP (both from F5).
\emph{Intuition.} The realized cost of immediacy without any quote data; widening $S^{\mathrm{eff}}$ predicts higher required returns (liquidity premium) and marks stress regimes in which reversal factors (F9--F11) pay more.

\factor{Expected time-to-next-trade}{Palguna--Pollak $\E[\tau\,|\,S]$}
\[
X_t = \widehat{\E}\big[\tau_t \,\big|\, S_t\big], \qquad \tau_t = \inf\{u>0 : \text{a print occurs at } t+u\},
\]
estimated nonparametrically over states $S_t$ (see F31) or with a simple hazard model on time-since-last-trade.
\emph{Intuition.} In a LOB this is the expected time to the first mid change; in bonds the expected time to the next print is itself a liquidity factor (staleness risk) and a gate: prediction horizons shorter than $\widehat{\E}[\tau]$ are not tradable.

\factor{Quote-update intensity}{toxic-flow feature (c)}
\[
U_t(\delta) = \#\{\text{composite-quote revisions in } [t-\delta, t]\}.
\]
\emph{Intuition.} Dealer quote-revision activity leads trade activity; a burst of revisions with no prints signals imminent repricing. One of the eight per-clock LOB statistics in the toxicity feature grid (F30).

%======================================================================
\section{Momentum and trend-reliability factors}\label{sec:mom}

\factor{Studentized (t-stat) momentum}{FactorMiner 023}
\[
X_t = \TsRank\!\left( \frac{\Delta(P,\,k)_t}{\Mean\big(|\Delta(P,\,k)|,\, n\big)_t},\; n\right).
\]
\emph{Intuition.} The trend divided by its own typical absolute size --- a signal-to-noise-scaled momentum that discounts moves within normal churn.
\emph{Bond translation.} Apply to OAS changes: 3--6 month spread momentum is a documented credit factor; the studentization suppresses stale-price artifacts.

\factor{Trend-reliability ($R^2$) weighting}{FactorMiner 080/083/085/090}
Regress price (or log price) on time over the window, obtain slope $\hat\beta_t$ and fit $R^2_t$; define
\[
X_t = \sgn(\hat\beta_t)\cdot R^2_t
\qquad\text{or}\qquad
X_t = \Rank(\hat\beta_t)\cdot \Rank(R^2_t).
\]
\emph{Intuition.} A trend that explains its own path is more likely to be information-driven and to continue; a low-$R^2$ ``trend'' is pressure and mean-reverts. Provides a data-driven split between the momentum and reversal families.

\factor{Cross-sectional lead--lag rank momentum}{FactorMiner 010/024/025}
\begin{gather*}
X_t = -\,\Rank\!\Big( \Rank\big(P_t/P_{t-k}\big) - \Rank\big(V_t/V_{t-k}\big) \Big),\\
X_t = \TsRank\big(\Delta(P^{\text{first}},k),n\big) - \TsRank\big(\Delta(P^{\text{last}},k),n\big),
\end{gather*}
where $P^{\text{first}}, P^{\text{last}}$ are first/last prints in each bar (open/close analogues).
\emph{Intuition.} Divergence between where the bar opens and closes, or between price and volume ranks, captures intrabar pressure direction.

%======================================================================
\section{Volatility and higher-moment factors}\label{sec:vol}

\factor{Realized volatility and inverse-vol tilt}{toxic-flow feature (a); FactorMiner 009/015}
\[
\sigma_t(n) = \Std(r,\,n)_t, \qquad
X_t = -\,\Rank(r_t)\cdot \Rank\big(\sigma_t(n)\big),
\qquad
X_t = -\,\Rank\big(\sigma_t(n)\big).
\]
\emph{Intuition.} (i) Mid-price volatility is a core state variable in every predictive grid in the corpus; (ii) reversal scaled up in high-vol names; (iii) the low-volatility tilt --- low-vol bonds earn higher risk-adjusted returns --- is a documented credit anomaly.

\factor{Volatility-of-volatility}{FactorMiner 077}
\[
X_t = \Std\big(\sigma_\cdot(n),\, m\big)_t .
\]
\emph{Intuition.} Instability of the volatility state flags regime transitions; in credit, vol-of-vol rises ahead of downgrades and liquidity spirals.

\factor{Rolling skewness and kurtosis regime switches}{FactorMiner 039/042--045/095}
\begin{gather*}
\Skew_t(n),\ \Kurt_t(n) \ \text{ of returns};\\
X_t = \operatorname{IfElse}\big( \Kurt_t(n) > \kappa,\; -\Rank(\Delta(P,k)),\; +\Rank(\Delta(P,k)) \big).
\end{gather*}
\emph{Intuition.} Higher moments as environment switches: fat-tailed regimes favor fading moves, thin-tailed regimes favor following them. Credit returns are structurally negatively skewed, so cross-sectional \emph{differences} in rolling skew are informative about crash exposure pricing. Estimation noise on sparse tapes is severe; use long windows and shrink toward bucket means.

%======================================================================
\section{Flow-toxicity and informed-flow factors}\label{sec:tox}

These originate in a broker's classification problem, but the \emph{label} (does the price move against the liquidity provider after the trade?) and the \emph{features} are pure prediction objects: a high predicted toxicity for observed flow is a directional signal that price will keep moving in the flow's direction.

\factor{Markout label and toxicity propensity}{Cartea--Duran-Martin--S\'anchez-Betancourt}
For a print at time $t$ with direction $d$, define the markout and label at horizon $G$:
\[
M_{t,G} = d\cdot\big(m_{t+G} - m_t\big), \qquad y_t = \ind{M_{t,G} > c} \in \{0,1\},
\]
and fit a classifier $\widehat{\Prob}(y_t = 1 \mid x_t)$ on the feature vector $x_t$ below. The \emph{factor} for bond $b$ is the toxicity-weighted signed flow
\[
X_t = \sum_{t_k \in [t-\delta, t]} \widehat{\Prob}(y_{t_k}=1 \mid x_{t_k})\; d_k V_k .
\]
\emph{Feature grid} (183 features in the source): eight LOB statistics --- (a) mid volatility, (b) client trade count, (c) quote-update count, (d) mid return, (e)--(f) $\log(1+q^B)$, $\log(1+q^A)$, (g) spread, (h) imbalance --- each under \emph{three clocks} (time, volume, transaction) and seven nested lookbacks, plus level features: order size, spread and imbalance at arrival, last bid/ask/mid, cumulative activity counts, a volatility estimate, and (crucially) the client's \emph{historical toxicity proportion}. Signed-log transforms $\sgn(x)\log(1+|x|)$ stabilize scale.
\emph{Bond translation.} Horizon $G$ of 1--5 days; mid replaced by the F15 fair value. Where client identity is unavailable (public TRACE), replace client features with trade-size buckets and side-of-market persistence; where it is available (dealer's own RFQ log), the historical-toxicity ratio is the strongest single feature.

\factor{Client historical toxicity ratio}{same}
\[
\rho^c_t = \frac{\#\{\text{toxic trades of client } c \text{ before } t\}}{\#\{\text{trades of client } c \text{ before } t\}} .
\]
Listed separately because it is a one-line feature with outsized importance, and because a pooled model over all clients \emph{with} client features outperforms per-client models --- a data-efficiency lesson that transfers to per-bond vs.\ pooled factor models.

\factor{VPIN (volume-synchronized probability of informed trading)}{Easley et al., via Li et al.}
Partition the tape into consecutive volume buckets of size $V^\ast$; within bucket $n$ let $V^{+}_n, V^{-}_n$ be buy and sell volume (by flag, or by bulk classification). Then
\[
\operatorname{VPIN}_t = \frac{1}{N}\sum_{n=1}^{N} \frac{|V^{+}_n - V^{-}_n|}{V^\ast}.
\]
\emph{Intuition.} A volume-clocked estimate of the informed share of flow; elevated VPIN precedes liquidity withdrawal and large moves. Its \emph{unsigned} nature makes it a risk/conditioning factor: it predicts volatility and the profitability of reversal factors rather than direction.

\factor{Informed-inventory signal (flow decomposition)}{Cartea--S\'anchez-Betancourt}
The dealing framework shows the relevant state variables are the informed traders' inventory and the two volume streams. As a prediction object: decompose observed flow using per-trade toxicity propensities (F27),
\begin{gather*}
Q^{\mathcal I}_t = \sum_{t_k \le t} \widehat{\Prob}(y_{t_k}=1\mid x_{t_k})\, d_k V_k,\\
Q^{\mathcal U}_t = \sum_{t_k \le t} \big(1-\widehat{\Prob}(y_{t_k}=1\mid x_{t_k})\big)\, d_k V_k,
\end{gather*}
and use as factors
\begin{gather*}
X_t = +\Rank\big(\Delta Q^{\mathcal I}_t\big) \quad (\text{informed flow: continuation}),\\
X_t = -\Rank\big(\Delta Q^{\mathcal U}_t\big) \quad (\text{uninformed flow: reversion}).
\end{gather*}
\emph{Intuition.} The two components of flow have opposite-signed predictive content; conflating them (as raw F3 does) averages the two effects away.

%======================================================================
\section{Intensity-ratio probability factors}\label{sec:ratio}

\factor{Next-trade-side probability (Cox intensity ratio)}{Muni Toke--Yoshida (both papers)}
Model buy and sell arrival intensities with a \emph{common unobserved baseline},
\[
\lambda^{\pm}(t) = \lambda_0(t)\, \exp\!\Big( \vartheta^{\pm}_0 + \sum_j \vartheta^{\pm}_j X_j(t) \Big),
\]
so that the baseline cancels in the ratio, leaving a logistic model with relative coefficients $\theta_j = \vartheta^{+}_j - \vartheta^{-}_j$:
\[
p_t \;=\; \Prob(\text{next trade is a buy} \mid \mathcal F_t)
\;=\; \frac{\lambda^{+}(t)}{\lambda^{+}(t)+\lambda^{-}(t)}
\;=\; \Big(1 + e^{-\theta_0 - \sum_j \theta_j X_j(t)}\Big)^{-1},
\qquad X_t = 2p_t - 1 .
\]
Covariates $X_j$: imbalance (F6/F5's $\Pi$), last sign (F1), signed spread (F2), Hawkes states (F4). Estimated by multinomial-logistic quasi-likelihood on event times only.
\emph{Intuition.} For bonds this construction is exceptionally well suited: the \emph{level} of activity is wildly nonstationary (new issues, index rebalances, macro prints) and lives in $\lambda_0$, which is never estimated; the tradable object --- the relative propensity of buys vs.\ sells given the state --- is exactly what survives the ratio. Two model-selection lessons transfer: imbalance alone is never the selected model, and the retained interaction is $s\epsilon$, not $s$.

\factor{Two-stage side $\times$ aggressiveness ratios}{marked-point-process extension}
Stage 1: side probability $p_t$ as in F31. Stage 2: conditional on side, the probability the trade is \emph{aggressive} (prints at or through the far composite quote), from a second ratio model with the same covariate family. The product gives four flow-type probabilities; the aggressive-buy minus aggressive-sell probability gap is the directional factor:
\[
X_t = p_t\, a^{+}_t - (1-p_t)\, a^{-}_t .
\]
\emph{Intuition.} Aggressive prints carry the information; the marked/two-stage decomposition isolates them. The same framework extends to predicting the \emph{size mark} of the next trade (volume prediction), useful as a flow forecast input to F13.

%======================================================================
\section{State-conditional nonparametric predictors}\label{sec:state}

\factor{Conditional moment grid over discretized states}{Palguna--Pollak}
Define a low-dimensional state, e.g.
\[
S_t = \Big( \big\lfloor \text{spread proxy} \big\rfloor,\ \big\lfloor \log(1+\text{time-since-last-trade}) \big\rfloor,\ \sgn(\text{last flow}) \Big),
\]
and estimate by historical averaging (with $k$-nearest-neighbor smoothing across adjacent states, respecting the sign of any imbalance coordinate):
\[
\widehat{\Prob}\big[\Delta m_{t,\delta} \neq 0 \,\big|\, S\big],\qquad
\widehat{\E}\big[\Delta m_{t,\delta} \,\big|\, S\big],\qquad
\widehat{\E}\big[\Delta m_{t,\tau} \,\big|\, S\big],\qquad
\widehat{\E}\big[\tau \,\big|\, S\big],
\]
where $\tau$ is the time of the first fair-value change (or next print, for bonds). Prediction is two-stage: first test ``will it move within $\delta$'' against a threshold, then take the sign of the conditional mean.
\emph{Intuition.} Fully nonparametric, hence assumption-free about functional form; the price is state coarseness. The two-stage structure (move / no-move, then direction) matches bond reality, where the modal outcome at short horizons is ``no print.''
\emph{Bond translation.} $\widehat{\E}[\tau\mid S]$ doubles as F23. Estimate within liquidity buckets, pooling bonds to populate states.

%======================================================================
\section{Composite and machine-generated cross-sectional alphas}\label{sec:mined}

Representative mined expressions from the RL/MCTS/agentic miners, translated to the bond tape (all crash-safe with $\varepsilon$ guards; $r$ denotes bar returns, $V$ bar volume):
{\small
\begin{align}
X^{(1)}_t &= -\,\Rank\!\left( \frac{P_t - \operatorname{Min}(P,n)_t}{\operatorname{Max}(P,n)_t - \operatorname{Min}(P,n)_t + \varepsilon} \right)
&&\text{(range-position fade)}\\[2pt]
X^{(2)}_t &= \Rank\big(\Delta(V,1)_t\big) - \Rank\!\Big(\tfrac{P_t - \VWAP_t}{\VWAP_t}\Big)
&&\text{(volume-lead divergence)}\\[2pt]
X^{(3)}_t &= -\,\TsRank\!\Big( K_t(n)\,,\ n \Big)\cdot \Rank\!\Big( \big|\Cov(r, V, n)_t\big| \Big)
&&\text{(range top $\times$ flow coupling)}\\[2pt]
X^{(4)}_t &= -\,\Rank\!\Big(\tfrac{P_t-\VWAP_t}{\VWAP_t}\Big)\cdot \big(1 - \Rank(V_t)\big)
&&\text{(low-volume dislocation fade)}\\[2pt]
X^{(5)}_t &= -\,\Rank\!\Big( \tfrac{\min(P^{\mathrm f}_t,\, P^{\mathrm l}_t) - \operatorname{Min}(P,n)_t}{R_t(n) + \varepsilon} \Big)
&&\text{(lower shadow, with $P^{\mathrm f},P^{\mathrm l}$ first/last prints)}\\[2pt]
X^{(6)}_t &= -\,\sgn\big(\Delta(P,1)_t\big)\cdot \TsRank\big(\Delta(P,k)_t,\, n\big)
&&\text{(sign-conditional fade)}
\end{align}}
Two structural lessons from the mining papers matter more than any single expression. First, alphas are evaluated and admitted \emph{as a pool}: the objective is the performance of the optimized linear combination, with mutual-correlation penalties, not individual information coefficients (Algorithm~\ref{alg:pool}). Second, the miners' libraries saturate: after $\sim$100 admitted factors, new candidates are near-duplicates --- diminishing returns are structural, so a bond implementation should expect a small number of genuinely distinct families (roughly the sections of this document).

%======================================================================
\section{Reproduction algorithms}\label{sec:algos}

All algorithms are stated language-agnostically; no code. Throughout, ``bar'' means an aggregation of the tape under a chosen clock (Algorithm~1).

\begin{algo}{Clock construction and fallback windows}
\item Choose a clock: time (fixed $\delta$), volume (fixed par amount $V^\ast$), or transaction (fixed count $k$). For bonds prefer volume or transaction clocks for flow features and time clocks for calendar-aligned fair values.
\item For each observation time $t$ and each nested horizon $\delta_1 < \dots < \delta_W$, collect the prints in the window ending at $t$.
\item If a window contains no prints, fall back to the next longer horizon; mark the observation missing only if the full-history window is empty.
\item Record, for every feature, which horizon actually supplied the data --- the realized horizon is itself a staleness feature.
\end{algo}

\begin{algo}{Decayed signed flow and impact state (F3, F13)}
\item Fix a half-life grid, e.g.\ $\{1, 5, 21\}$ trading days; convert to decay rates $\beta = \log 2 / \text{half-life}$.
\item Iterate over prints in time order, updating $F \leftarrow F e^{-\beta \Delta t} + d_k V_k / \text{AmtOutstanding}$ at each print.
\item For the impact state, estimate $(\lambda, \beta)$ per liquidity bucket: regress the future $h$-day return on $F_t(\beta)$ for each candidate $\beta$ over a grid of horizons $h$; pick the $\beta$ maximizing out-of-sample fit, and set $\lambda$ to the (negated) regression slope rescaled by $(1-e^{-\beta h})$.
\item Output the factor $X_t = -\lambda F_t(\beta)$ and the de-pressured value $P_t - \lambda F_t(\beta)$.
\end{algo}

\begin{algo}{Markout labeling and toxicity propensity (F27--F28)}
\item Choose a markout horizon $G$ (1--5 days) and a materiality threshold $c$ (e.g.\ half the typical effective spread).
\item For every print, compute the fair value at $t$ and at $t+G$ (F15 median of prints, or composite mid where available); label $y=1$ if the fair value moved in the trade's direction by more than $c$.
\item Build the feature grid: the eight statistics (vol, trade count, quote-update count, return, log bid depth, log ask depth, spread, imbalance) under each clock and each nested horizon; append level features and, if identity is known, the client's historical toxicity ratio; apply $\sgn(x)\log(1+|x|)$ to signed magnitudes and $\log(1+x)$ to volumes.
\item Fit any calibrated classifier with strictly walk-forward training (the label needs $G$ of future data --- embargo accordingly).
\item Score every new print; aggregate propensity-weighted signed volume over a trailing window to obtain the directional factor, and split flow into informed/uninformed components for F30.
\end{algo}

\begin{algo}{Intensity-ratio (next-side) estimation (F31)}\label{alg:ratio}
\item Collect the event times of customer buys and sells for the estimation universe (pool sparse bonds within buckets).
\item At each event, record the covariate vector (imbalance or volume pressure, last sign, signed spread, Hawkes states, decayed flow).
\item Maximize the multinomial-logistic quasi-likelihood: each event contributes $\log$ of the model probability of the side that occurred; equivalently, fit a logistic regression of the event side on the covariates sampled \emph{at event times}.
\item Select covariates with an information criterion (BIC-type penalties preferred; AIC over-selects here) or $\ell_1$ penalty; expect interactions ($s\epsilon$) to survive and single-covariate models to be rejected.
\item Output $X_t = 2\hat p_t - 1$ evaluated continuously between events.
\end{algo}

\begin{algo}{Sparse fair-value feature selection (Section \ref{sec:fv})}
\item Extract the full grid: percentiles, min, max, first, last, mean, VWAP, volatility, counts, volumes --- per side and combined, per nested horizon (with fallback), yielding a few hundred candidates.
\item Standardize features cross-sectionally within date and bucket.
\item Fit $\ell_1$-penalized regression (or quantile regression if a distributional forecast is wanted) of the $h$-day forward return on the grid, choosing the penalty by walk-forward validation.
\item Read off the retained features and their importance by summed absolute coefficients across type, horizon, and side; expect percentile/min/max levels at medium horizons to dominate and volume features to be dropped.
\end{algo}

\begin{algo}{State-conditional moment tables (F33)}
\item Discretize the state (spread bucket, staleness bucket, last-flow sign) coarsely enough that each cell has hundreds of observations after pooling within liquidity buckets.
\item Sweep the historical tape once, accumulating per-cell counts and sums for: move indicator at horizon $\delta$, move size, move size at first change, and time to first change.
\item At prediction time, average the cell estimates over the $k$ nearest states, treating states with opposite imbalance/flow sign as infinitely distant.
\item Predict in two stages: move vs.\ no-move against a threshold on the conditional probability, then the sign of the conditional mean.
\end{algo}

\begin{algo}{Pool-level factor combination (mining-style)}\label{alg:pool}
\item Maintain a factor pool with weights. For each candidate factor, tentatively add it and re-optimize the weights of the linear combination against the target (rank-IC of the combined score with forward returns).
\item Admit the candidate only if the \emph{pool} objective improves net of a redundancy penalty (e.g.\ cap pairwise factor correlation, or penalize the increase in combination-weight collinearity).
\item Periodically prune: remove any factor whose deletion does not degrade the pool objective.
\item Freeze weights on a schedule; all evaluation strictly walk-forward, cross-sectionally within buckets.
\end{algo}

\begin{algo}{Liquid-to-illiquid propagation (F18)}
\item Map every bond to its issuer curve; designate the most liquid sibling (highest trailing print count) as the anchor.
\item Compute all factors on the anchor; assign them to the illiquid siblings, optionally decayed by curve distance (maturity gap).
\item Additionally compute the sibling's basis to the anchor-implied curve value; use the basis as a convergence factor.
\item Never propagate in the reverse direction (illiquid-estimated models onto liquid bonds); the generalization is empirically asymmetric.
\end{algo}

%======================================================================
\section{Source map}\label{sec:map}

\begin{center}
\small
\begin{tabular}{@{}p{5.6cm}lp{6.2cm}@{}}
\toprule
Source (as read) & Tag & Factors drawn \\
\midrule
Lehalle \& Neuman, \emph{Incorporating signals into optimal trading} & LN & F6 (imbalance as OU signal) \\
Muni Toke \& Yoshida, \emph{Ratios of Cox-type intensities} & MTY20 & F1, F2, F31 \\
Muni Toke \& Yoshida, \emph{Marked point processes and intensity ratios} & MTY22 & F4, F31, F32 \\
Palguna \& Pollak, \emph{Mid-price prediction in a LOB} & PP16 & F8, F23, F33; share-imbalance variant of F6 \\
Li et al., \emph{Intelligent market making} & LDZ14 & F7; VPIN pointer (F29) \\
Cartea, Duran-Martin \& S\'anchez-Betancourt, \emph{Detecting toxic flow} & CDS & F24 (quote-update count), F27, F28; clock grid \\
Cartea \& S\'anchez-Betancourt, \emph{A simple strategy to deal with toxic flow} & CSB & F30 (informed/uninformed decomposition) \\
Bank, Cartea \& K\"orber, \emph{Optimal execution and speculation with trade signals} & BCK & F3 (flow-driven price changes; liquidity state) \\
Webster (chapter), \emph{Mathematical models of price impact} & OW/W & F13 (propagator, myopia of impact--alpha link) \\
Yu et al., \emph{OrderFusion} & OF & F5, F22 (side separation) \\
Yu et al., \emph{Orderbook feature learning} (electricity) & OFL & F15--F18; Algorithm 5; asymmetric generalization \\
Yu, Xue, Ao et al., \emph{AlphaGen} (KDD'23) & AG & operator grammar; Algorithm 7 \\
Ren et al., \emph{RiskMiner} & RM & operand/operator inventory \\
Wang et al., \emph{FactorMiner} & FM & F9--F12, F14, F19--F21, F25--F26, F34--F36 grid, Section \ref{sec:mined} expressions \\
Yu, Xue, Ao \& He, \emph{LLM-assisted alpha discovery} & LLM & grammar validation conventions \\
\bottomrule
\end{tabular}
\end{center}

\paragraph{Exclusions.} Quote-placement and dealing-strategy mathematics (optimal spreads, inventory controls, HJB systems), execution-scheduling solutions, and broker internalization/externalization policies are out of scope; only the predictive state variables inside those frameworks were retained.

\end{document}
